\chapter{冒険心のつまみ}
% full version: chapters/08_nora.tex

どこで昼食をとるか選ぶとしよう。行きつけのカフェに行くか、新しい店を試すか。いつも行きつけばかりなら、角を曲がった先にもっといい店があっても一生知らずに終わる。いつも新しい店ばかり試していれば、まずい昼食に当たることも多い。ちょうどよい中間はどこにあるのか。探すか、見つけたものを使うか。この問題は、学習するあらゆる生き物の前にある。脳でこれを担っているのは、どうやら放送局\nt{NORA}、つまりノルアドレナリンらしい。

\section*{コントラストのつまみ}

ノルアドレナリンは、ニューロンが信号にどれだけ鋭く反応するかを変える。作用が弱いと、ニューロンはなめらかな調光器に似ている。信号が少し増えれば応答も少し明るくなる。作用が強いと、スイッチに似てくる。閾値より下では沈黙、上では全力の応答。より正確に言えば、実験によると、ノルアドレナリンはニューロンの背景のおしゃべりを、本物の信号への応答より強く抑える。「コントラストのつまみ」は、この効果を扱いやすくしたモデルである。

\pencilfig{8}{\pencilart{8}}{冒険心のつまみ：左に回すとマシンは新しいものを試し、右に回すと行きつけを迷わず選ぶ。}

\section*{考えるか、決めるか}

第3章の、互いに抑制し合う2つの群れを思い出そう。コントラストが低いと両方が働き、弱い方が少し静かになる。ネットワークはまだ「考えて」いて、選択肢を比べている。コントラストがある閾値を超えると、ネットワークは急に「決めた」モードに移る。一方の群れが勝ち、もう一方は黙る。1つのつまみが、マシンを熟考と決断の間で切り替えるのだ。

\section*{どれだけ探すか}

ノイズを加えよう。コントラストが低いと、ノイズは選択肢どうしの小さな差を簡単に上回り、マシンは最善でない選択肢をたびたび試す。つまり探す。コントラストが高いと、ノイズは何も左右せず、マシンはいつも最善と思うものを選ぶ。つまり使う。

私たちのモデルでは、世界はときどき変わる。最善の選択肢が最悪になるのだ。成果はつまみの位置に対して逆U字形になる。コントラストが小さすぎると、マシンは手探りでさまよう。大きすぎると、世界が変わったことに気づかず、まずくなったカフェに意地を張って通い続ける。そして世界が頻繁に変わるほど、最適な設定は低くなる。変わりやすい世界では、探す方が得なのだ。

\pencilfig{108}{%
% points: models/nora_explore.py, reward as a share of the best possible, gain 0.5 ... 64 (doubling); the y axis starts at 0.70, hence the break mark
\draw[pen] (0,0) -- (6.3,0); \draw[pen] (0,0) -- (0,0.18); \draw[pen] (0,0.34) -- (0,2.4);
\draw[pcGray, line width=0.6pt] (-0.1,0.14) -- (0.1,0.22); \draw[pcGray, line width=0.6pt] (-0.1,0.3) -- (0.1,0.38);
\node[lbl, anchor=north] at (3.0,-0.05) {コントラストのつまみ};
\node[lbl, anchor=south, rotate=90] at (-0.1,1.3) {成果};
\draw[penlite=pcBlue] plot[smooth] coordinates {(0.00,0.55) (0.85,0.79) (1.70,1.19) (2.55,1.68) (3.40,2.00) (4.25,2.07) (5.10,2.01) (5.95,1.91)};
\draw[penlite=pcRed] plot[smooth] coordinates {(0.00,0.55) (0.85,0.69) (1.70,0.93) (2.55,1.24) (3.40,1.40) (4.25,1.28) (5.10,1.06) (5.95,0.89)};
\draw[dotpen=pcBlue, wash=pcBlue] (4.25,2.07) circle (0.07);
\draw[dotpen=pcRed, wash=pcRed] (3.40,1.40) circle (0.07);
\node[lbl, text=pcBlue, anchor=south] at (4.25,2.15) {世界がまれに変わる};
\node[lbl, text=pcRed, anchor=north] at (3.40,1.30) {頻繁に};
\node[lbl, anchor=north west, align=left] at (0.05,-0.35) {手探りで\\さまよう}; \node[lbl, anchor=north east, align=right] at (6.3,-0.35) {変化に\\気づかない};
}{つまみの位置ごとのモデルの成果。どの曲線にも頂点があり、世界が頻繁に変わるほど頂点は低く、左に寄る。}

\section*{つまみを回すのは誰か}

変化の自然な兆候は予想外であること、つまり誤差がいつもより大きくなることだ。ある仮説によれば、ノルアドレナリンが知らせているのはまさにこうした予想外の変化である。ここで微妙な点が大事になる。一定の背景レベルのノルアドレナリンはむしろ探索や注意散漫と結びつき、重要な出来事に対する短い連打は集中と決断に結びつく。連打はコンピュータの割り込みのように働くのかもしれない。「今の作業をやめろ、世界が変わった、計画を立て直せ」。

このシステムの間接的な兆候は外からも見える。何かが予想外に変わり、人がより速く学び始めると、瞳孔が広がる。ただし瞳孔はノルアドレナリンだけを反映しているわけではないので、これは手がかりであって計測器ではない。

そして正直に告白しておく。私たちのモデルでは、予想外さに応じてつまみを調整しても、最良の固定設定に比べてほんのわずかしか得にならなかった。こうした調整がどこで本当に元をとるのかは、これから解明すべきことである。

\fullversion{8}
